\bookchapter{Git: Best Practices and Recovery}\label{ch:git4}

\section*{Explanation}
Knowing the commands is half of it. The other half is the habits that keep a repository readable and a team fast, and the calm to fix mistakes. Everybody makes mistakes with Git: committing to the wrong branch, a typo in a message, a deleted file. Almost all of them are recoverable, if you know which tool fits. This chapter first collects the habits of good teams, then the repair kit.

\section*{Commits}
\subsection*{Atomic commits}
A commit should contain \emph{one logical change} that builds and works. ``Fix the camera timeout'' is one commit. ``Fix the camera, rename variables, update README and try a new gait'' is four. Small commits are easy to review, to understand months later, to undo (\texttt{git revert}) and to find bugs in (\texttt{git bisect}, below).

If you have mixed several changes in one file, stage them piece by piece:
\begin{verbatim}
$ git add -p walk.py
@@ -10,3 +10,4 @@ def step(speed):
 ...
Stage this hunk [y,n,q,a,d,s,e,?]? y
\end{verbatim}
\texttt{y} stages this piece, \texttt{n} skips it, \texttt{s} splits it into smaller pieces.

\subsection*{Good commit messages}
The log is the diary of the project. A message has a short subject, an empty line and, when needed, a body:
\begin{verbatim}
Fix camera timeout after Wi-Fi reconnect

The camera node needs up to 4 s to resume after a reconnect, but
the controller stopped after 2 s. Raise the timeout to 5 s and log
a warning when more than 2 s are used.

Closes #12
\end{verbatim}
Rules that most teams follow:
\begin{itemize}[nosep]
  \item Subject: at most about $50$ characters, no full stop, \textbf{imperative mood} (``Fix'', ``Add'', ``Remove'', as if giving an order to the code). It should complete the sentence ``If applied, this commit will\dots''
  \item Empty second line. Body lines at most $72$ characters.
  \item The body explains \textbf{why}, not what: the diff already shows what changed.
  \item Refer to issues (\texttt{Closes \#12}).
\end{itemize}
\begin{center}
\renewcommand{\arraystretch}{1.25}
\begin{tabular}{p{6.0cm}p{7.2cm}}
\textbf{Poor} & \textbf{Good}\\ \hline
\texttt{fix} & \texttt{Fix crash when camera returns None}\\
\texttt{changes} & \texttt{Move gait parameters to config.yaml}\\
\texttt{wip, again, final2} & \texttt{Add logging of loop time per cycle}\\
\texttt{Updated stuff in walk.py} & \texttt{Reduce step height to avoid tripping}\\
\end{tabular}
\end{center}
Without \texttt{-m}, \texttt{git commit} opens your editor for a long message. Do not commit generated files, secrets or debug leftovers (see Chapter~\ref{ch:git1}); run the tests first.

\section*{Branching strategies}
\begin{description}
  \item[Trunk-based / GitHub flow (recommended).] One long-lived branch \texttt{main} that always works. Everything else is a short branch ($\leq$ few days), merged by PR (Chapter~\ref{ch:git3}). Releases are tags on \texttt{main}. Simple, and enough for almost every student or robotics team.
  \item[Git flow.] Additional long-lived branches \texttt{develop} and \texttt{release/*}. More ceremony; it makes sense for software with several supported versions installed at customers.
\end{description}
Whatever the choice, write it down in the \texttt{README} or \texttt{CONTRIBUTING.md}, so that everybody does the same: branch names, how to merge (squash or rebase), who reviews.

\subsection*{Your daily routine}
\begin{center}
\renewcommand{\arraystretch}{1.25}
\begin{tabular}{p{2.6cm}p{10.4cm}}
\textbf{When} & \textbf{What}\\ \hline
Morning & \verb!git switch main && git pull! (start from the newest state)\\
New task & \verb!git switch -c feature/name!\\
While working & small commits; \verb!git status!, \verb!git diff --staged! before each; run the tests\\
Often & \verb!git fetch! and \verb!git rebase origin/main! (or merge) to stay close to \texttt{main}\\
Before the PR & \verb!git log --oneline main..! (is the story clean?), \verb!git rebase -i! if needed, push\\
Evening & push your branch (a backup!), even if unfinished\\
After merge & switch to \texttt{main}, pull, delete the branch\\
\end{tabular}
\end{center}
Settings that make the routine smoother:
\begin{verbatim}
$ git config --global pull.rebase true       # no needless merge commits
$ git config --global fetch.prune true       # drop deleted remote branches
$ git config --global rerere.enabled true    # remember conflict solutions
$ git config --global alias.lg "log --oneline --graph --decorate --all"
$ git lg
\end{verbatim}

\section*{Undoing things}
The right tool depends on \emph{how far the mistake has gone}. Figure~\ref{fig:g4_undo} is the decision, Table~\ref{tab:g4_undo} lists the commands.
\begin{center}
\begin{tikzpicture}[q/.style={draw,rounded corners,thick,fill=yellow!20,minimum width=4cm,minimum height=0.9cm,align=center,font=\footnotesize},a/.style={draw,rounded corners,thick,fill=green!12,minimum width=4.6cm,minimum height=0.9cm,align=center,font=\footnotesize\ttfamily},>={Stealth[length=2.5mm]}]
  \node[q] (q1) at (0,0) {Where is the mistake?};
  \node[q] (q2) at (0,-1.7) {Not yet committed};
  \node[q] (q3) at (0,-3.4) {Committed, not pushed};
  \node[q] (q4) at (0,-5.1) {Pushed / shared};
  \draw[->,thick] (q1)--(q2); \draw[->,thick] (q2)--(q3); \draw[->,thick] (q3)--(q4);
  \node[a] (a2) at (6.6,-1.7) {git restore FILE};
  \node[a] (a3) at (6.6,-3.4) {git commit --amend\\git reset};
  \node[a] (a4) at (6.6,-5.1) {git revert COMMIT};
  \draw[->,thick] (q2)--(a2); \draw[->,thick] (q3)--(a3); \draw[->,thick] (q4)--(a4);
  \node[font=\scriptsize,align=left,gray!60!black,anchor=west] at (9.0,-1.7) {discard edits};
  \node[font=\scriptsize,align=left,gray!60!black,anchor=west] at (9.0,-3.4) {rewrite my local history};
  \node[font=\scriptsize,align=left,gray!60!black,anchor=west] at (9.0,-5.1) {new commit that cancels};
\end{tikzpicture}
\captionof{figure}{Which undo tool? The more people have seen the commit, the less history you may rewrite.}\label{fig:g4_undo}
\end{center}
\begin{center}
\captionof{table}{The undo commands.}\label{tab:g4_undo}
\renewcommand{\arraystretch}{1.25}
\begin{tabular}{p{5.6cm}p{7.6cm}}
\textbf{Command} & \textbf{Effect}\\ \hline
\verb!git restore FILE! & throw away changes in the working directory (\textbf{cannot be undone})\\
\verb!git restore --staged FILE! & unstage, keep the edits\\
\verb!git restore --source=HEAD~2 FILE! & bring an old version of one file back\\
\verb!git commit --amend! & change the last commit (message or content)\\
\verb!git reset --soft HEAD~1! & undo the commit, keep changes staged\\
\verb!git reset HEAD~1! (mixed) & undo the commit, keep changes unstaged\\
\verb!git reset --hard HEAD~1! & undo the commit \emph{and} delete the changes\\
\verb!git revert COMMIT! & new commit that undoes an old one; safe for shared history\\
\verb!git reflog! & the diary of where HEAD has been: the safety net\\
\end{tabular}
\end{center}

\begin{verbatim}
$ git commit -m "Update REAMDE"        # typo in the message
$ git commit --amend -m "Update README"
$ echo broken > walk.py                # disaster in the working dir
$ git restore walk.py                  # gone: back to the last commit
\end{verbatim}
\textbf{Amend and reset only for commits that you have not pushed}, as they replace commits by new ones (the golden rule of Chapter~\ref{ch:git2}). For a pushed commit use \texttt{revert}:
\begin{verbatim}
$ git revert 4dc63e4
[main 4809596] Revert "Add a"
 1 file changed, 1 deletion(-)
\end{verbatim}

\subsection*{The safety net: reflog}
You ran \verb!git reset --hard HEAD~2! and the two newest commits are gone from the log! But Git remembers every place HEAD has been for about $90$ days:
\begin{verbatim}
$ git reflog
a36b8f4 HEAD@{0}: reset: moving to HEAD~2
7698a27 HEAD@{1}: commit: Add b
4dc63e4 HEAD@{2}: commit: Add a
a36b8f4 HEAD@{3}: commit (amend): Update README
$ git reset --hard 7698a27         # or HEAD@{1}: back to before
HEAD is now at 7698a27 Add b
\end{verbatim}
Commits are not erased immediately; only the labels moved. This is why \emph{committed work is almost never lost}, while \emph{uncommitted work is}: \texttt{git restore} and \texttt{reset --hard} cannot bring back edits that were never committed. Commit often, also unfinished work on your own branch.

\subsection*{Force-push}
After a rebase or amend of a branch that you already pushed, a normal push is rejected, since the histories differ. On \emph{your own feature branch} that nobody else uses you may overwrite it:
\begin{verbatim}
$ git push --force-with-lease
\end{verbatim}
\texttt{--force-with-lease} refuses if somebody pushed in the meantime, so it cannot silently destroy their work, unlike plain \texttt{--force}. On \texttt{main} protect against force-pushes in the branch rules (Chapter~\ref{ch:git3}).

\section*{Finding the commit that broke it: bisect}
The robot walked fine two weeks ago and now falls over, with $64$ commits in between. \texttt{git bisect} does a binary search through the history:
\begin{verbatim}
$ git bisect start
$ git bisect bad                   # the current version is broken
$ git bisect good v0.1             # this tag was fine
Bisecting: 31 revisions left to test (roughly 5 steps)
$ # test now (run the program), then tell git:
$ git bisect good                  # or: git bisect bad
...
8cc1a09 is the first bad commit
$ git bisect reset                 # back to where you started
\end{verbatim}
Since $2^6=64$, about six tests find the culprit among $64$ commits. With a test script (\verb!git bisect run ./test.sh!) it runs by itself. This only works well with atomic commits that each build.

\section*{Large files and Git LFS}
Git stores every version of every file forever. A $200$\,MB model committed ten times makes the repository $2$\,GB for every clone. Keep datasets and weights out of Git (a storage server, release assets, or a tool such as DVC). If binaries must live next to the code, \textbf{Git LFS} (Large File Storage) keeps them on a separate server, and the repository holds only small pointers:
\begin{verbatim}
$ sudo apt install git-lfs && git lfs install
$ git lfs track "*.onnx"           # writes .gitattributes
$ git add .gitattributes policy.onnx
\end{verbatim}
GitHub rejects single files above $100$\,MB without LFS.

\section*{Hooks and automatic checks}
A \textbf{hook} is a script that runs at a Git event. A \texttt{pre-commit} hook can run a formatter and a quick test before each commit, so errors never enter the history. The \texttt{pre-commit} tool manages them from a file \texttt{.pre-commit-config.yaml}:
\begin{verbatim}
$ pip install pre-commit        # in your virtual environment
$ pre-commit install            # activate for this repository
\end{verbatim}
The same checks should also run in GitHub Actions (Chapter~\ref{ch:git3}), since a hook can be skipped on a single laptop but the pull request check cannot.

\section*{Working across Windows, macOS and Linux}
Windows ends text lines with two characters (CRLF), Linux with one (LF). A shell script saved with CRLF fails on the robot with the confusing message \verb!bad interpreter: /bin/bash^M!. Let Git handle it: a file \texttt{.gitattributes} in the repository with the line \verb!* text=auto eol=lf! stores LF in the repository, and Windows users set \verb!git config --global core.autocrlf true!. Also avoid file names that differ only in upper and lower case, since Windows treats them as the same.

\section*{Checklist}
\begin{keyidea}[Ten habits]
(1)~\texttt{git status} and \texttt{git diff --staged} before every commit. (2)~One logical change per commit. (3)~Messages that explain why. (4)~A branch for every task, merged by PR. (5)~Short-lived branches. (6)~\texttt{main} always works; tests run automatically. (7)~Pull before you start, push before you leave. (8)~Never rewrite shared history. (9)~No secrets, no large binaries, a good \texttt{.gitignore}. (10)~When in doubt: \texttt{git status}, then \texttt{git reflog}, and do not panic.
\end{keyidea}

\begin{example}
Ada wanted to throw away her last (unpushed) commit but typed \verb!git reset --hard HEAD~3! instead of \verb!HEAD~1!. Three commits seem to be gone. She runs \texttt{git reflog}, finds the line \texttt{HEAD@\{1\}: commit: Add logging} that was her newest commit, and runs \verb!git reset --hard HEAD@{1}!. Everything is back. Had she also had uncommitted edits, the \texttt{--hard} would have destroyed them, and the reflog could not help: this is why she commits before experimenting.
\end{example}

\section*{Practice problems}
\begin{enumerate}[label=\textbf{\thechapter.\arabic*}]
  \item Rewrite the commit message ``fixed stuff'' for a change that stops the controller from crashing when the IMU sends \texttt{NaN}. Obey the usual rules.
  \item Why should commits be small and atomic? Give two benefits.
  \item You edited \texttt{config.py} for two unrelated reasons. How do you create two separate commits?
  \item You made a typo in the message of your last, unpushed commit. Which command fixes it? Why would you not do this for a pushed commit?
  \item Which command discards all uncommitted changes in \texttt{walk.py}? Can you get them back?
  \item Explain the difference between \verb!git reset --soft HEAD~1!, \verb!git reset HEAD~1! and \verb!git reset --hard HEAD~1!.
  \item A faulty commit is already on \texttt{main} and your colleagues have pulled it. How do you undo its effect?
  \item You ran \verb!git reset --hard! and lost commits. How do you find them again?
  \item Why is \texttt{--force-with-lease} better than \texttt{--force}? When is force-pushing acceptable?
  \item[\textbf{\thechapter.10}$\star$] Between the tag \texttt{v1.0} (works) and \texttt{main} (broken) there are $100$ commits. About how many tests does \texttt{git bisect} need at most? Write the commands to start it.
\end{enumerate}
